\mode*
\section{Related work}
\label{relatedwork}

With the catalogue and its analyses in place, we can now situate our work among
three strands of literature: applications of variation theory outside
introductory programming, teaching informed by students' misconceptions but
without variation theory, and concrete proposals for how to teach introductory
programming.

\subsection{Variation theory beyond programming}

Variation theory has a long tradition as a design theory for teaching in other
subjects \autocite{MartonPang2006Conditions}. In the \emph{learning-study}
tradition --- lesson study with
variation theory as its explicit learning theory --- the theory serves as a
source of guiding principles for teachers' pedagogical design, lesson analysis
and evaluation \autocite{PangLing2011LearningStudy}. Mathematics education is
the most developed case: using variation and invariance within and between
instructional examples helps learners engage with mathematical structure
\autocite{Kullberg2017VTMathematics}. In science education,
\textcite{Preston2022ElectricityVT} combine variation theory with multimodal
representations to design a teaching sequence on electric circuits for primary
school, and \textcite{Kullberg2024PlanningTeaching} collect case studies of
teaching informed by variation theory.

Within computing education, the phenomenographic tradition that our method
builds on goes back to \textcite{Booth1992Phenomenographic}, who studied how
students experience learning to program. Closest to our work,
\textcite{EckerdalThune2005ObjectsVT} study first-year students' understandings
of \emph{object} and \emph{class} phenomenographically and connect them to
variation theory, and \textcite{ThuneEckerdal2009ConceptionsVT} apply
phenomenography and variation theory to students' conceptions of computer
programming as such, identifying dimensions of variation to use in teaching.
These works apply variation theory to a single concept pair or to programming
as a whole; we instead started from the misconceptions documented in the
literature and derived, for each of them (\cref{misconceptions}), the critical
aspects and a pattern of variation to teach against it.

\subsection{Teaching from misconceptions, without variation theory}

Using students' conceptions as the starting point for instruction is the core
of the \emph{conceptual change} tradition in science education, going back to
\textcite{Posner1982ConceptualChange}; contemporary models emphasise that
conceptual change is dynamic and iterative, with regressions and multiple entry
points \autocite{Nadelson2018ConceptualChange}. The conceptions themselves are
elicited with diagnostic instruments --- most commonly interviews, open-ended
tests, multiple-choice tests and multiple-tier tests
\autocite{KaltakciGurel2015Diagnostics} --- of which concept inventories are
the best-known form. \Textcite{Julie2020ConceptInventory} develop such an
inventory for programming, covering variables, conditionals and functions,
explicitly to show teachers their students' mental models.

In computing education, \textcite{Ma2011MentalModels} investigate and improve
novices' non-viable mental models of value and reference assignment;
\textcite{Sorva2013NotionalMachines} reviews the misconception, mental-model
and threshold-concept literatures and argues that an explicitly taught
\emph{notional machine} addresses novices' difficulties with the runtime
dynamics of programs; and \textcite{Gusukuma2018MDF} generate feedback from
detected misconceptions rather than from output checks alone.
\Textcite{MisconceptionsSurvey2017} survey this literature. These approaches
share our premise that misconceptions should drive teaching, but they derive
different things from it: diagnosis, feedback, or an explicit model of the
machine. Variation theory adds a prescription for the \emph{content} of
instruction: which aspect must vary, against which invariant background, for
the critical aspect to become discernible.

\subsection{Concrete approaches to teaching introductory programming}

A third strand proposes concrete teaching designs for introductory programming
without deriving them from specific misconceptions. At course level,
\textcite{Medeiros2018SLR} review the challenges of teaching and learning
introductory programming and categorise them, \textcite{Schulte2006TeachersTeach}
chart what teachers actually teach and what their students find difficult, and
\textcite{Miliszewska2007Befriending} propose a balanced course design in
response to novice difficulties. At the level of classroom activities,
well-evaluated designs include subgoal-labelled worked examples, which reduced
withdrawal and failure in introductory programming
\autocite{Margulieux2020Subgoals}; Parsons problems, which practise basic
programming constructs in a puzzle-like format
\autocite{Parsons2006Puzzles}; peer instruction, which halved fail
rates in computing courses \autocite{Porter2013PeerInstruction}; and PRIMM,
which structures lessons as predict--run--investigate--modify--make
\autocite{Sentance2019PRIMM}.

These designs specify the \emph{activity structure} --- what students do ---
whereas our analyses specify the \emph{example content}: which examples to
juxtapose so that the critical aspect of a documented misconception comes into
focus. The two are complementary; the predict step of PRIMM and the questions
of peer instruction are natural carriers for the contrast patterns we derived
in the catalogue chapters.
